% GENERATED FILE. Edit canonical lessons, then run npm run generate:books.
% canonical-source-hash: fnv1a64:34e61bba8f458faa
% canonical-lessons: FR-C245-semblable, FR-C245-juste, FR-C245-faux-fausse, FR-C245-pret-prete, FR-C245-gratuit-gratuite

\chapter{Similar, Fair, False, Ready, Free of charge}
\label{ch:fr-similar-fair-false-ready-free-of-charge}
\hlchaptermodality{245}

\begin{chapteropening}
\textbf{By the end of this chapter:} \emph{I can say similar, fair, false, ready and free of charge.}

Complete the last lesson of chapter 245: I can say similar, fair, false, ready and free of charge.
\end{chapteropening}

\section[semblable]{semblable — similar}

\label{lesson:FR-C245-semblable}



\begin{quote}
Before the new one: say the French for smooth, then the French for rough.
\end{quote}

\subsection*{You'll want to know: semblable}
\textbf{semblable} — \textquotedblleft{}similar\textquotedblright{}.

\begin{grammarlens}[title={What you've built}]
One more word for this chapter.
\end{grammarlens}

\subsection*{Your turn}
\begin{itemize}
\raggedright
  \item \emph{Say it:} \emph{semblable}
  \item \emph{Say it:} \emph{semblable}, once more
  \item \emph{Say it:} say \emph{semblable}
  \item \emph{Recall:} say the French for smooth, then the French for rough, then say \emph{semblable} again
\end{itemize}

\begin{tcolorbox}[breakable,colback=teal!4,colframe=teal!35!black,title={Before you move on}]
What does semblable mean? (\textquotedblleft{}Similar\textquotedblright{}.) Say it once more.
\end{tcolorbox}
\section[juste]{juste — fair, right}

\label{lesson:FR-C245-juste}



\begin{quote}
Before the new one: say the French for rough, then the French for similar.
\end{quote}

\subsection*{You'll want to know: juste}
\textbf{juste} — \textquotedblleft{}fair, right\textquotedblright{}.

\begin{grammarlens}[title={What you've built}]
One more word for this chapter.
\end{grammarlens}

\subsection*{Your turn}
\begin{itemize}
\raggedright
  \item \emph{Say it:} \emph{juste}
  \item \emph{Say it:} \emph{juste}, once more
  \item \emph{Say it:} say \emph{juste}
  \item \emph{Recall:} say the French for rough, then the French for similar, then say \emph{juste} again
\end{itemize}

\begin{tcolorbox}[breakable,colback=teal!4,colframe=teal!35!black,title={Before you move on}]
What does juste mean? (\textquotedblleft{}Fair, right\textquotedblright{}.) Say it once more.
\end{tcolorbox}
\section[faux / fausse]{faux / fausse — false, wrong}

\label{lesson:FR-C245-faux-fausse}



\begin{quote}
Before the new one: say the French for similar, then the French for fair.
\end{quote}

\subsection*{You'll want to know: faux / fausse}
\textbf{faux / fausse} — \textquotedblleft{}false, wrong\textquotedblright{}.

\begin{grammarlens}[title={What you've built}]
One more word for this chapter.
\end{grammarlens}

\subsection*{Your turn}
\begin{itemize}
\raggedright
  \item \emph{Say it:} \emph{faux / fausse}
  \item \emph{Say it:} \emph{faux / fausse}, once more
  \item \emph{Say it:} say \emph{faux / fausse}
  \item \emph{Recall:} say the French for similar, then the French for fair, then say \emph{faux / fausse} again
\end{itemize}

\begin{tcolorbox}[breakable,colback=teal!4,colframe=teal!35!black,title={Before you move on}]
What does faux / fausse mean? (\textquotedblleft{}False, wrong\textquotedblright{}.) Say it once more.
\end{tcolorbox}
\section[prêt / prête]{prêt / prête — ready}

\label{lesson:FR-C245-pret-prete}



\begin{quote}
Before the new one: say the French for fair, then the French for false.
\end{quote}

\subsection*{You'll want to know: prêt / prête}
\textbf{prêt / prête} — \textquotedblleft{}ready\textquotedblright{}.

\begin{grammarlens}[title={What you've built}]
One more word for this chapter.
\end{grammarlens}

\subsection*{Your turn}
\begin{itemize}
\raggedright
  \item \emph{Say it:} \emph{prêt / prête}
  \item \emph{Say it:} \emph{prêt / prête}, once more
  \item \emph{Say it:} say \emph{prêt / prête}
  \item \emph{Recall:} say the French for fair, then the French for false, then say \emph{prêt / prête} again
\end{itemize}

\begin{tcolorbox}[breakable,colback=teal!4,colframe=teal!35!black,title={Before you move on}]
What does prêt / prête mean? (\textquotedblleft{}Ready\textquotedblright{}.) Say it once more.
\end{tcolorbox}
\section[gratuit / gratuite]{gratuit / gratuite — free of charge}

\label{lesson:FR-C245-gratuit-gratuite}



\begin{quote}
Before the new one: say the French for false, then the French for ready.
\end{quote}

\subsection*{You'll want to know: gratuit / gratuite}
\textbf{gratuit / gratuite} — \textquotedblleft{}free of charge\textquotedblright{}.

\begin{grammarlens}[title={What you've built}]
One more word for this chapter.
\end{grammarlens}

\subsection*{Your turn}
\begin{itemize}
\raggedright
  \item \emph{Say it:} \emph{gratuit / gratuite}
  \item \emph{Say it:} \emph{gratuit / gratuite}, once more
  \item \emph{Say it:} say \emph{gratuit / gratuite}
  \item \emph{Recall:} say the French for false, then the French for ready, then say \emph{gratuit / gratuite} again
\end{itemize}

\begin{tcolorbox}[breakable,colback=teal!4,colframe=teal!35!black,title={Before you move on}]
What does gratuit / gratuite mean? (\textquotedblleft{}Free of charge\textquotedblright{}.) Say it once more.
\end{tcolorbox}
